% Figure: the permutation null distribution of a difference in group
% means. The histogram is the reference distribution built by relabelling
% the pooled data; the vertical line is the observed difference. The
% bar heights are precomputed counts from a representative run of the
% worked example in the text (B = 2000 relabellings, binned).
\begin{figure}[htbp]
\centering
\begin{tikzpicture}
\begin{axis}[
    width=12cm, height=6cm,
    ybar interval=0.9,
    xmin=-9, xmax=9,
    ymin=0, ymax=0.20,
    axis lines=left,
    xlabel={Permuted difference in means $\bar{x}_A^{\star} - \bar{x}_B^{\star}$ (working units)},
    ylabel={Relative frequency},
    xtick={-8,-6,-4,-2,0,2,4,6,8},
    ytick={0,0.05,0.10,0.15,0.20},
    every axis plot/.append style={fill=engblue!30, draw=engblue!60!black},
]
% Approximately symmetric, mean zero, sd about 2.4 working units.
\addplot+[hist] coordinates {
  (-9,0.003) (-8,0.006) (-7,0.013) (-6,0.027) (-5,0.049)
  (-4,0.081) (-3,0.116) (-2,0.146) (-1,0.158) (0,0.146)
  (1,0.116) (2,0.081) (3,0.049) (4,0.027) (5,0.013)
  (6,0.006) (7,0.003) (8,0.000)
};
% Observed statistic at +5.4
\draw[engred, very thick] (axis cs:5.4,0) -- (axis cs:5.4,0.18);
\node[engred, anchor=south, font=\small] at (axis cs:5.4,0.18) {observed $5.4$};
\end{axis}
\end{tikzpicture}
\caption[Permutation null distribution of a difference in means]{The
permutation null distribution of the difference in group means,
built by pooling the two groups and recomputing the statistic under
$B = 2000$ random relabellings. The reference distribution is centred
on zero, the value the statistic would take if the group label
carried no information. The observed difference of $5.4$ working
units sits in the upper tail; the two-sided permutation $p$-value is
the fraction of relabelled statistics at least as extreme in
absolute value as the observed one. The construction assumes only
exchangeability of the observations under the null, not a Gaussian
model, which is what makes it the engineer's default when the
distributional assumption of the $t$-test is in doubt.}
\label{fig:vol02:ch14:permutation-null}
\end{figure}
